\begin{proposition}[Fundamental group of an irreducible fibre]\hypertarget{tag-02Y2}{}\label{night:P4R.1}\label{stmt:P4R.1}
\StatementRecord{night:P4R.1}{proved}{02Y2}{the definitions and constructions stated in the proof}

For a proper holomorphic degeneration of complex two-tori with
smooth total threefold and central fibre $mS$, where $S$ is
irreducible and reduced, the group $\pi_1(S)$ has an abelian
normal subgroup with cyclic quotient of order dividing $m$.
\end{proposition}
\begin{proof}
Write the degeneration as $f:X\to\Delta$.
The compatible triangulation of {\L}ojasiewicz
\cite[Theorem~2, p.~464]{Lojasiewicz1964} applies to
the smooth real analytic total space and its closed analytic
subset $S$. Its hypotheses hold because the total space is
countable and finite dimensional, and the collection consisting
of $S$ is locally finite and semianalytic. Let $\tau:|K|\to X$ be the resulting homeomorphism and
$L\subset K$ the subcomplex with $|L|=\tau^{-1}(S)$.
After barycentric subdivision, $L$ is a full subcomplex.
For barycentric coordinates $\lambda_v$, put
$q=\sum_{v\in L}\lambda_v$. On the open star $q>0$, the map
\[
 r_K\Bigl(\sum_v\lambda_v v\Bigr)
   =\sum_{v\in L}\frac{\lambda_v}{q}v
\]
erases the other coordinates and divides the remaining ones
by their sum. For $0\le s\le1$, the interpolation
$H_s(x)=(1-s)x+s r_K(x)$ stays in the same star, since
$q(H_s(x))=(1-s)q(x)+s>0$, and fixes $L$.
Transport by $\tau$ gives a strong deformation
retraction $r:V\to S$. Compactness of $S$ and successive
subdivision make $V$ lie inside a relatively compact
neighbourhood of $S$.

Properness gives a small disc $\Delta_\epsilon$ with
$U=f^{-1}(\Delta_\epsilon)\subset V$: on the compact
complement of $V$ in a larger closed tube, $|f|$ has a
positive minimum. Thus the inclusion $i:S\to U$ and the
continuous restriction $r:U\to S$ satisfy $r\circ i=\id$.
The neighbourhood-retraction statement immediately after
Definition~1 of Fern\'andez de Bobadilla--Koll\'ar
\cite[p.~85]{BobadillaKollar2012} applies to the same proper
map $f:X\to\Delta$.

Write $U^*=U\setminus S$. The proper smooth fibration
$U^*\to\Delta_\epsilon^*$ has fibre $T^4$. Its homotopy
sequence and $\pi_2(\Delta_\epsilon^*)=0$ give
\[
 \pi_1(U^*)=\Lambda\rtimes\langle t\rangle,
 \qquad \Lambda=\mathbb Z^4.
\]
Choose $t$ mapping to $1$ in the infinite cyclic quotient.
Every element is uniquely $\lambda t^n$, with $\lambda\in\Lambda$;
conjugation by $t$ gives the displayed semidirect product.
The inclusion $U^*\to U$ induces a surjection on fundamental
groups. A loop can be perturbed off the real-codimension-two
smooth stratum of $S$ and off its higher-codimension singular
strata. The same general-position argument for a null-homotopy
disc avoids the singular strata and meets $S_{\rm reg}$ in
finitely many transverse points. Deleting little discs about
these intersections expresses the kernel as the normal closure
of the corresponding meridians. The smooth locus of the
irreducible $S$ is connected. Along a path in that locus, a
trivialization of its complex normal line gives an annulus of
positive meridians in $U^*$, proving their conjugacy. The complex
normal coordinate fixes their positive orientation.

At a smooth point of $S$, write $f=vz^m$ with $v$ a
holomorphic unit. On a small coordinate neighbourhood,
$z'=z\exp((\log v)/m)$ is a holomorphic normal coordinate
and $f=(z')^m$. For small $a>0$, the positive circle
$z'=a\exp(2\pi i\theta)$, $0\le\theta\le1$, maps to
$a^m\exp(2\pi im\theta)$, so its image under
$\pi_1(U^*)\to\mathbb Z$ is $m$. Let $A$ be the image of $\Lambda$
in $\pi_1(U)$. It is abelian and normal, and
$\pi_1(U)/A$ is cyclic with order dividing $m$.
The homomorphism $r_*:\pi_1(U)\to\pi_1(S)$ is surjective
because $r_*i_*=\id$. Its image $r_*(A)$ is abelian and
normal, and the induced quotient remains cyclic of order
dividing $m$. This proves the assertion for $\pi_1(S)$.

\end{proof}